
The Software Defined Networking (SDN) paradigm \cite{Nunes2014ASO} enables to decouple the control plane from the forwarding plane of network devices. This separation allows their dynamic configuration using interfaces such as OpenFlow \cite{OpenFlow}. A recent development, based on P4 \cite{P4}, offers to go further and enables program deployment into the devices in addition to configuration. Several studies investigate the use of P4 for detection of heavy network flows (which may interfere with smaller flows sharing the same path and thus impact the QoS) or for traffic detection. 
\cite{P4_HeavyHitter} addresses the detection of large flows with an algorithm consisting of a hash table pipeline, which identifies 95\% of the heaviest flows on an ISP trace. Other works show threshold-based techniques, either by distributing thresholding across all the switches \cite{P4_HeavyHitter2} or by focusing on their size and duration \cite{P4_HeavyHitter3}. Some works focus on flow classification using ML models on P4 switches. Inter-arrival times (IATs) histograms are used by \cite{P4_Seior} to determine which TCP flavour (Cubic or BBR) controls each flow and to classify each flow on P4 hardware using a ML model on the control plane.  
\cite{P4_Switches_ML, P4_Flow_Classif} directly implement ML models on the programmable switches to handle traffic classification.
\cite{P4_Switches_ML} maps four trained ML models (DT, K-means, SVM, and Naïve Bayes) to match-action pipelines and achieve full-line rate classification.
\cite{P4_Flow_Classif} proposes \textit{SMASH-D1} (resp. \textit{SMASH-D2}) to reduce the DT model complexity for flow classification by performing the classification decision one (resp. two) level(s) above the DT leaf nodes.

However, many research papers implement their solution with the P4 Linux software switch (BMv2), and very few on a real P4 hardware switch. This is a crucial point since the  software implementation is largely less restrictive than the hardware one. Indeed, a hardware switch must process packets at line rate and has less computational facilities (e.g. no floating-point operations and limited memory-accesses.)  than a Linux system. \cite{P4_Switches_ML} outlines that by adapting their approach to run on hardware-based prototypes, classifiers can use fewer features and classes.  \cite{9838674} investigated to what extent a P4 hardware solution for AQM can be achievable and faced limitations imposed by the P4 hardware. They then split the processing between the P4 data plane for simple operations and control plane for complex computations. 

Our proposal differs from the related work in several ways: (i) unlike \cite{P4_HeavyHitter, P4_HeavyHitter3, P4_Flow_Classif}, we perform the traffic detection using a P4 hardware implementation instead of a software one; (ii) we split the CG traffic detection task by performing the features extraction in the data plane with P4 switches and the CG traffic detection in the control plane with a NFV node; (iii) unlike the previous works \cite{NOMS}, we use unsupervised learning for CG traffic detection to cope with the generalization problem we may face with supervised learning \cite{LearningClassif}.

